%%=============================================================================
%% Conclusie
%%=============================================================================

\chapter{Conclusie}%
\label{ch:conclusie}

Deze bachelorproef onderzocht de centrale onderzoeksvraag uit sectie~\ref{sec:onderzoeksvraag}: \textit{``Hoe effectief is gamification bij het aanleren van COBOL aan junior developers die gewend zijn aan moderne talen zoals Python?''}
Om hierop een onderbouwd antwoord te formuleren, werd een literatuurstudie uitgevoerd (hoofdstuk~\ref{ch:stand-van-zaken}), een Proof-of-Concept ontwikkeld als interactieve webapplicatie (hoofdstuk~\ref{ch:poc}), en de PoC geëvalueerd in twee delen (hoofdstuk~\ref{ch:evaluatie}). Eerst via een uitgevoerde literatuurevaluatie. Daarna via een voorbereid maar niet uitgevoerd pilotontwerp voor een vergelijking met traditionele methoden zoals een papieren cursus.

\section{Antwoord op de onderzoeksvragen}

\begin{description}
    \item[DV1: uitdagingen in syntax en logica.] Inhoudelijk beantwoord in hoofdstuk~\ref{ch:stand-van-zaken} (sectie~\ref{sec:COBOL-versus-python}). Belangrijke verschilpunten voor een Python-georiënteerde junior zijn onder meer de \texttt{PIC}-clausule, de hiërarchische level-numbers, de scope-valstrik met de punt (\texttt{.}) en de verbose programmastructuur.

    \item[DV2: barrières bij het leren.] Beschreven in hoofdstuk~\ref{ch:stand-van-zaken}. Bovenop de syntactische verschillen geldt de \emph{didactische kloof}: het bestaande COBOL-leermateriaal is vooral statisch, terwijl de doelgroep gewend is aan interactieve sandboxes.

    \item[DV3: gamification-elementen.] De literatuur onderscheidt een bredere set gamification-\allowbreak elementen dan enkel de PBL-triade: prestatiebeloningen (punten, badges, mijlpalen), uitdagingselementen (levels en opdrachten), sociale elementen (leaderboards en vergelijking), narratieve/rol-\allowbreak elementen (thema, verhaal, avatars), en feedback- en progressie-\allowbreak elementen (onmiddellijke validatie, checklists, voortgangsbalken) \autocite{KevinWerbach2020, Chou2015}. Voor het leren van COBOL sluiten vooral drie categorie\"en het best aan: (1) uitdagingen met snelle feedback, (2) zichtbare prestatiebeloningen, en (3) herhaalbare oefening in een levelstructuur (distributed practice/spaced repetition). Die combinatie ondersteunt zowel \textit{Development \& Accomplishment} als \textit{Empowerment of Creativity \& Feedback} en past bij de specifieke leerdrempels van COBOL (strikte structuur, expliciete syntax en veel procedurele patronen), terwijl sociale competitie via leaderboards best beperkt en doordacht wordt ingezet omdat dit niet voor elke beginner motiverend werkt \autocite{Zhan2022, KevinWerbach2020, Cepeda2006}.

    \item[DV4: ontwerp van een gamified leeromgeving.] Concreet uitgewerkt in hoofdstuk~\ref{ch:poc}. De literatuurevaluatie in hoofdstuk~\ref{ch:evaluatie} linkt de belangrijkste ontwerpkeuzes aan onderzoek (punten en badges, leaderboard, Octalysis, MDA en didactische feedback). Daardoor is het ontwerp als geheel aannemelijk op basis van de geraadpleegde literatuur (secties~\ref{sec:evaluatie-pbl} tot~\ref{sec:evaluatie-didactisch}). De beperkingen daarvan staan expliciet in sectie~\ref{sec:evaluatie-beperkingen}.

    \item[DV5: leer-effectiviteit t.o.v.\ traditioneel.] Op basis van de literatuurevaluatie is het aannemelijk dat de gamification-elementen in \textit{Learn COBOL} een positief effect hebben op motivatie en leerwinst, in dezelfde grootteorde als de meta-analyse van \textcite{Zhan2022} rapporteert voor gamification in programmeeronderwijs (\(SMD \approx 0{,}77\) op motivatie, \(\approx 0{,}75\) op prestaties). Een rechtstreekse vergelijking met traditionele methoden zoals een papieren cursus is in deze bachelorproef niet getoetst; een definitief, kwantitatief antwoord op DV5 vraagt metingen met deelnemers. In sectie~\ref{sec:evaluatie-pilotstudie} staat daarvoor een pilotontwerp dat de leeruitkomsten van twee groepen vergelijkt, één die COBOL leert via de PoC en één die COBOL leert via een papieren cursus. Die pilot is in deze bachelorproef niet uitgevoerd.
\end{description}

\section{Bijdrage}

De toegevoegde waarde van dit onderzoek is viervoudig: een theoretisch kader dat gamification (PBL, MDA en Octalysis) koppelt aan mainframe-onderwijs; een functionele Proof-of-Concept op moderne webtechnologie (React, TypeScript, Supabase, Monaco) uitgewerkt in hoofdstuk~\ref{ch:poc}; een uitgevoerde literatuurevaluatie die per ontwerpkeuze onderbouwt waarom het ontwerp aannemelijk is; en een herbruikbaar pilotontwerp (sectie~\ref{sec:evaluatie-pilotstudie}) dat de stap naar empirische validatie concreet maakt.

\section{Aanbevelingen voor empirische opvolging}
\label{sec:conclusie-empirische-opvolging}

De voornaamste aanbeveling is de in hoofdstuk~\ref{ch:evaluatie} beschreven pilot daadwerkelijk uit te voeren. Dat is geen vervanging van de literatuurevaluatie, maar de volgende stap om de verwachtingen te toetsen met deelnemers en om voorzichtiger te kunnen spreken over oorzaak en overdraagbaarheid.
Het onderzoeksdesign staat al volledig uitgewerkt in sectie~\ref{sec:evaluatie-pilotstudie}: doelstellingen, doelgroep, een tussen-proefpersoon opzet met willekeurige toewijzing aan groep A (papieren cursus) of groep B (PoC), een tijdslijn van vijf weken, en een gedeelde kennistest aan het eind van de leerdag als hoofdmeting.
Wat in deze bachelorproef ontbreekt, is enkel de uitvoering met deelnemers.